\chapter{Reproducibility and Resource Integration}
\label{chap:reproducibility}

This chapter documents how the supplied resources are integrated into the thesis. The repository contains a SaarTex thesis template, multiple manuscript versions, workshop material, a poster, figures, and bibliography files. The rewritten thesis does not treat these resources as unrelated fragments. It synthesizes them into one coherent document about adversarial robustness of SSL vision encoders across downstream tasks.

\section{Resource Categories}

The SaarTex template provides the document class, title structure, declarations, and bibliography setup. The thesis keeps this structure because it is the correct local format. The main manuscript resources provide the abstract, introduction, background, method, results, conclusion, appendix, macros, and bibliography for the SSL robustness study. The ICML workshop and poster resources provide concise versions of the same motivation, methods, figures, and results. The figure resources provide the attack-interface diagram, \DeACL illustration, and fine-tuning curves.

A good rewrite must use these resources without simply pasting them verbatim. Conference papers are short and optimized for page limits. A thesis must be more explanatory. Therefore, the rewritten chapters expand definitions, explain design choices, interpret tables slowly, and separate related work from methodology. The numbers remain grounded in the supplied results, while the prose is expanded into thesis form.

\section{Template Integration}

The main document uses `Thesis.tex` as the root. It declares the title, author, thesis subject, location, advisor placeholders, theorem-like math macros, and custom commands for attack and model names. The content chapters are included in order. This structure keeps front matter, main matter, bibliography, and appendix separate.

The front matter includes the abstract, table of contents, and sworn declaration. The main matter now includes introduction, background, related work, methodology, experimental design, evaluation, extended analysis, reproducibility, and discussion. The appendix includes technical details. This expanded structure is closer to a thesis than the previous compact draft.

The `Options.tex` file sets typography, bibliography style, and packages. The caption package is included because the evaluation chapter uses unnumbered captions inside a multi-panel figure layout. The graphics path points to the `Figures/` directory, where the supplied result and framework figures are available.

\section{Bibliography Integration}

The bibliography combines SSL, adversarial robustness, dense prediction, dataset, and depth-estimation references. It includes foundational SSL works such as SimCLR, SimSiam, \DINO, MAE, and \DINOv \parencite{chen2020simple,simsiam,caron2021dino,mae,oquab2024dinov2}. It includes adversarial robustness foundations and robust SSL methods \parencite{biggio2013evasion,szegedy2014intriguing,madry2018towards,kim2020adversarial,jiang2020robust,fan2021does,zhang2022decoupled,luo2023rethinking}. It also includes segmentation and depth references \parencite{long2017segmentation,xiao2018unified,godard2019depthestimation,MegaDepthLi18,AdaBins}.

The rewrite avoids introducing new claims that require unsupported references. Where a concept is explained generally, it is tied to existing citations. Where a number is reported, it comes from the supplied result tables and poster material. This keeps the document internally consistent.

\section{Figure Integration}

The attack-interface figure is used in the methodology chapter. It explains that an SSL encoder can be attacked through its embedding space or through downstream heads. This visual is important because the thesis argument is structural: robustness depends on the interface. A figure communicates that structure more efficiently than prose alone.

The \DeACL figure is used in the defense section. It shows robust fine-tuning as a transfer from a pre-trained encoder to a robust encoder. The figure supports the explanation of the two-part objective: preserve clean teacher representations and align adversarial examples with clean robust representations.

The classification, segmentation, and depth fine-tuning curves are used together in the evaluation chapter. They show that \DeACL fine-tuning dynamics differ by task and attack. The curves are not merely decoration; they support the claim that \EmbedAttack gains appear more clearly than task-specific robustness gains.

\section{Table Integration}

The result tables are central. They are not moved to the appendix because the thesis argument depends on them. Classification, segmentation, and depth each receive a table with clean, \EmbedAttack, and task-specific attack columns. This repeated structure makes the cross-task pattern easy to see.

The hyperparameter table is placed in the appendix. This keeps the main chapters readable while still recording enough detail for reproducibility. Hyperparameters include perturbation budget, attack steps, step size, optimizer choices, batch sizes, epochs, weight decay, and \DeACL fine-tuning settings.

The choice of table placement reflects thesis priorities. Main text should present evidence and interpretation. Appendix should preserve technical detail. A reader can understand the argument without memorizing every optimizer setting, but the settings remain available.

\section{Why the Case-Study Chapter Was Removed}

The previous draft included a trustworthy-AI case-study chapter about watermark localization, member-vs-generated inference, and federated-gradient reconstruction. Those topics are interesting, but they do not belong in a thesis that is supposed to use the SSL robustness resources as its core. Keeping them made the thesis broader but less coherent. It also made the title and abstract misleading.

The rewritten document therefore removes that chapter from the main thesis. This is not because provenance or privacy are unimportant. It is because the supplied resource set for the thesis topic is about adversarial robustness of self-supervised vision encoders across downstream tasks. A focused thesis should develop that topic deeply rather than mix in unrelated competitions.

If those case studies were required by an examiner, they should become a separate part with a different title, motivation, and contribution statement. Under the current instruction to use the resources correctly, the focused SSL robustness thesis is the cleaner choice.

\section{Length Expansion Strategy}

A conference paper compresses motivation, related work, methods, and results into a few pages. A thesis should unpack them. The expansion strategy used here has four parts. First, the introduction explains the research question and contributions in more detail. Second, a related-work chapter separates SSL, dense prediction, adversarial examples, and robust SSL. Third, an experimental-design chapter explains why each task, dataset, attack, and metric was chosen. Fourth, an extended-analysis chapter interprets the result tables beyond the short evaluation summary.

This expansion increases page count without inventing new experiments. It turns compact resource material into thesis material. That is important because padding a thesis with irrelevant content would be worse than being short. The added pages support the same central argument from different angles: deployment interface determines robustness.

\section{Reproducibility Limits}

The repository in this environment does not provide a complete runnable training and evaluation pipeline. It provides LaTeX resources, figures, and bibliography material. Therefore, reproducibility in this rewrite means document reproducibility and result traceability, not rerunning the full experiments. The numerical results are reproduced from supplied tables and figures.

A fully reproducible research artifact would also include training scripts, exact model checkpoints, dataset preprocessing, attack implementations, random seeds, hardware information, and logs. The thesis can describe these requirements, but it cannot create them from LaTeX-only resources. This limitation should be stated rather than hidden.

\section{Document Build Considerations}

The thesis uses standard LaTeX constructs: chapters, sections, equations, figures, tables, and BibLaTeX citations. A complete local build should run LaTeX, Biber, and LaTeX enough times to resolve citations and references. In this execution environment, TeX binaries may not be installed, so static checks are used as a fallback.

Static checks cannot replace compilation. They can catch missing citation keys, obvious brace imbalance, and stale references to removed content. A final human workflow should still compile the PDF locally, inspect figure placement, check overfull boxes, and verify the page count. This is especially important after expanding the thesis to target at least fifty pages.

\section{Expected Page Count After Expansion}

The expanded thesis is designed to pass the fifty-page target by increasing substantive source text and by adding chapters that introduce natural page breaks. The added related-work, experimental-design, extended-analysis, and reproducibility chapters contribute significant prose. Figures, tables, front matter, bibliography, and appendix add additional rendered pages.

The exact page count depends on the local LaTeX installation, fonts, float placement, bibliography rendering, and table placement. Nevertheless, based on source length, chapter count, and float count, the expanded document should be in the range expected for a fifty-page thesis. If a compiled PDF still falls short, the best next expansion would be to add a detailed per-dataset qualitative discussion and a limitations chapter with threat-model variants.

\section{Checklist for Final Review}

A final review should check five items. First, the title, abstract, and contributions should all describe the same thesis. Second, every chapter should support the SSL robustness question. Third, every table number should match the supplied resources. Fourth, no deleted case-study references should remain in the main document. Fifth, the bibliography should compile without missing keys.

The rewrite is structured to satisfy those items. It focuses the thesis, expands the explanation, and keeps the resource-derived numerical results visible. The remaining work after compilation is editorial: polish wording, adjust float placement, and fill in advisor/referee metadata.
